% Options for packages loaded elsewhere
\PassOptionsToPackage{unicode}{hyperref}
\PassOptionsToPackage{hyphens}{url}
\documentclass[
  10pt,
  a4paper,
]{article}
\usepackage{xcolor}
\usepackage[margin=22mm]{geometry}
\usepackage{amsmath,amssymb}
\setcounter{secnumdepth}{-\maxdimen} % remove section numbering
\usepackage{iftex}
\ifPDFTeX
  \usepackage[T1]{fontenc}
  \usepackage[utf8]{inputenc}
  \usepackage{textcomp} % provide euro and other symbols
\else % if luatex or xetex
  \usepackage{unicode-math} % this also loads fontspec
  \defaultfontfeatures{Scale=MatchLowercase}
  \defaultfontfeatures[\rmfamily]{Ligatures=TeX,Scale=1}
\fi
\usepackage{lmodern}
\ifPDFTeX\else
  % xetex/luatex font selection
\fi
% Use upquote if available, for straight quotes in verbatim environments
\IfFileExists{upquote.sty}{\usepackage{upquote}}{}
\IfFileExists{microtype.sty}{% use microtype if available
  \usepackage[]{microtype}
  \UseMicrotypeSet[protrusion]{basicmath} % disable protrusion for tt fonts
}{}
\makeatletter
\@ifundefined{KOMAClassName}{% if non-KOMA class
  \IfFileExists{parskip.sty}{%
    \usepackage{parskip}
  }{% else
    \setlength{\parindent}{0pt}
    \setlength{\parskip}{6pt plus 2pt minus 1pt}}
}{% if KOMA class
  \KOMAoptions{parskip=half}}
\makeatother
\setlength{\emergencystretch}{3em} % prevent overfull lines
\providecommand{\tightlist}{%
  \setlength{\itemsep}{0pt}\setlength{\parskip}{0pt}}
\usepackage{mathrsfs}
\newcommand{\Log}{\operatorname{Log}}
\usepackage{mathtools}
\usepackage{xurl}
\emergencystretch=3em
\setcounter{tocdepth}{1}
\newcommand{\Beta}{\mathrm{B}}
\DeclareUnicodeCharacter{00B2}{\ensuremath{^2}}
\DeclareUnicodeCharacter{00B1}{\ensuremath{\pm}}
\DeclareUnicodeCharacter{00D7}{\ensuremath{\times}}
\DeclareUnicodeCharacter{2192}{\ensuremath{\to}}
\DeclareUnicodeCharacter{21D2}{\ensuremath{\Rightarrow}}
\DeclareUnicodeCharacter{21A6}{\ensuremath{\mapsto}}
\DeclareUnicodeCharacter{2208}{\ensuremath{\in}}
\DeclareUnicodeCharacter{2209}{\ensuremath{\notin}}
\DeclareUnicodeCharacter{2212}{\ensuremath{-}}
\DeclareUnicodeCharacter{2227}{\ensuremath{\wedge}}
\DeclareUnicodeCharacter{2228}{\ensuremath{\vee}}
\DeclareUnicodeCharacter{2260}{\ensuremath{\ne}}
\DeclareUnicodeCharacter{2264}{\ensuremath{\le}}
\DeclareUnicodeCharacter{2265}{\ensuremath{\ge}}
\DeclareUnicodeCharacter{2297}{\ensuremath{\otimes}}
\DeclareUnicodeCharacter{00B3}{\ensuremath{^3}}
\DeclareUnicodeCharacter{2074}{\ensuremath{^4}}
\DeclareUnicodeCharacter{2076}{\ensuremath{^6}}
\DeclareUnicodeCharacter{03BB}{\ensuremath{\lambda}}
\DeclareUnicodeCharacter{03BC}{\ensuremath{\mu}}
\DeclareUnicodeCharacter{03B5}{\ensuremath{\epsilon}}
\DeclareUnicodeCharacter{03C4}{\ensuremath{\tau}}
\DeclareUnicodeCharacter{03A6}{\ensuremath{\Phi}}
\DeclareUnicodeCharacter{03B1}{\ensuremath{\alpha}}
\DeclareUnicodeCharacter{03C3}{\ensuremath{\sigma}}
\DeclareUnicodeCharacter{03B4}{\ensuremath{\delta}}
\DeclareUnicodeCharacter{03C0}{\ensuremath{\pi}}
\DeclareUnicodeCharacter{039B}{\ensuremath{\Lambda}}
\DeclareUnicodeCharacter{2202}{\ensuremath{\partial}}
\DeclareUnicodeCharacter{0393}{\ensuremath{\Gamma}}
\DeclareUnicodeCharacter{2205}{\ensuremath{\varnothing}}
\DeclareUnicodeCharacter{221E}{\ensuremath{\infty}}
\DeclareUnicodeCharacter{2200}{\ensuremath{\forall}}
\DeclareUnicodeCharacter{00B9}{\ensuremath{^1}}
\DeclareUnicodeCharacter{00B0}{\ensuremath{{}^\circ}}
\DeclareUnicodeCharacter{211D}{\ensuremath{\mathbb{R}}}
\usepackage{bookmark}
\IfFileExists{xurl.sty}{\usepackage{xurl}}{} % add URL line breaks if available
\urlstyle{same}
\hypersetup{
  hidelinks,
  pdfcreator={LaTeX via pandoc}}

\author{}
\date{}

\begin{document}

{
\setcounter{tocdepth}{3}
\tableofcontents
}
\section{The original endpoint pair: exact transport and the full Weil
matrix}\label{the-original-endpoint-pair-exact-transport-and-the-full-weil-matrix}

The two endpoint arguments of the original zeta theta identity supply a
pair of tests with no common nontrivial zero. We calculate their exact
time transport, the original signed divisor at every finite cutoff, and
the full translated Weil matrix. The scalar sign of one test does not
determine this matrix.

The incoming proofs are ER1--40a,
\href{supporting_proofs/endpoint_resonance_20260923/ENDPOINT_RESONANCE_AND_PRIME_CLASS.tex}{complete
ER source}, and AP1--28,
\href{supporting_proofs/arithmetic_pairing_20260923/ACTUAL_ENDPOINT_PAIRING.tex}{complete
AP source}. Their complete unchanged sources and exact hashes are
retained in the accompanying intake. ER and AP use the original
\href{https://arxiv.org/abs/1801.05914v5}{Rodgers--Tao author formulas,
arXiv:1801.05914v5}, equations phidef, htdef, hoz, and sas; the
explicit-formula convention is
\href{https://arxiv.org/abs/math/9811068v1}{Connes,
arXiv:math/9811068v1}, Appendix II (1)--(11). Root read both incoming
proofs completely. This note does not claim a new reading of either
author archive.

\subsection{1. Original function and endpoint
arguments}\label{original-function-and-endpoint-arguments}

Throughout \(t>0\), \(x,u\in\mathbb R\), \(s\in\mathbb C\), and
\(w=s-\tfrac12\). Retain \[
\begin{aligned}
\Phi(u)&=\sum_{n\ge1}(2\pi^2n^4e^{9u}-3\pi n^2e^{5u})e^{-\pi n^2e^{4u}},\\
H_t(x)&=\int_0^\infty e^{tu^2}\Phi(u)\cos(xu)\,du,\\
A(s)&=\pi^{-s/2}\Gamma(s/2),\quad P(w)=w^2-\tfrac14=s(s-1),\\
\mathcal Q(s)&=s(s-1)A(s)\zeta(s),\qquad
16H_0(-2i(s-\tfrac12))=\mathcal Q(s).
\end{aligned}\tag{EPM1}
\] The displayed product, including each exceptional factor, is retained
whenever \(\mathcal Q\) is used. It is an auxiliary entire product, not
the working replacement for \(\zeta\). In particular the original theta
identity is \[
A(s)\zeta(s)=-\frac1s+\frac1{s-1}
+\int_1^\infty\sum_{n\ge1}e^{-\pi n^2v}
\{v^{s/2-1}+v^{(1-s)/2-1}\}\,dv .
\tag{EPM2}
\] The integral is entire by exponential domination. Thus
\(\mathcal Q(0)=\mathcal Q(1)=1\). At \(s=-2m+h\), \(m\ge1\), the
original factors have the full local forms \[
\begin{aligned}
\zeta(-2m+h)&=h\,u_m(h),\quad
\mathcal Q(-2m+h)=b_m(h)u_m(h),\\
b_m(h)&=\frac{(-1)^m4m(2m+1)\pi^m}{m!}
(1-h/(2m))(1-h/(2m+1))\pi^{-h/2}\Gamma(1+h/2)
\prod_{\ell=1}^{m}(1-h/(2\ell))^{-1}.
\end{aligned}\tag{EPM3}
\] Here \(u_m,b_m\) are nonvanishing holomorphic germs. Gamma recurrence
proves every factor in \(b_m\); the original functional equation and the
nonzero Euler product at \(1+2m\) prove the simple zero and
\(u_m(0)\ne0\). At \(s=1+h\), retain \[
\zeta(1+h)=h^{-1}u_{\rm pole}(h),\quad
s(s-1)A(s)=h(1+h)\pi^{-(1+h)/2}\Gamma((1+h)/2),
\quad u_{\rm pole}(0)=1.
\tag{EPM4}
\] Consequently the zeros of \(\mathcal Q\) are exactly the original
nontrivial zeros, with the same multiplicities. They lie in
\(0<\Re s<1\). To include the boundary argument, absolute logarithmic
Euler expansion for \(\sigma>1\) and
\(3+4\cos v+\cos(2v)=2(1+\cos v)^2\ge0\) give \[
\zeta(\sigma)^3|\zeta(\sigma+iy)|^4|\zeta(\sigma+2iy)|\ge1.
\tag{EPM5}
\] If \(y\ne0\) and \(\zeta(1+iy)\) vanished to order \(m\ge1\), the
left side would be \(O((\sigma-1)^{4m-3})\), tending to zero. The other
nonreal factor is bounded because the only pole is at one. This is a
contradiction. The point one has the retained product value EPM4.
Reflection in EPM2 gives the zero-free opposite boundary. The Euler
product and the same reflection exclude the two exterior half-planes.

Every term of \(\Phi(u)\) is positive for \(u\ge0\), since
\(2\pi n^2e^{4u}>3\). Its even extension and all derivatives have
faster-than-Gaussian decay. Direct differentiation of the convergent
series and Fourier integration by parts show that \(H_t\) is real even
Schwartz on the real line for every real \(t\).

Define two real sources with their parity retained: \[
\begin{aligned}
h_{c,t}(x)&=e^{-x^2/(4t)}H_t(x)\cos(x/(2t)),&
h_{s,t}(x)&=e^{-x^2/(4t)}H_t(x)\sin(x/(2t)),\\
T&=D_x^2-\tfrac14,& f_{i,t}&=Th_{i,t}\quad(i=c,s),\\
M_i(t,s)&=\int_{\mathbb R}h_{i,t}(x)e^{-wx}\,dx,&
F_i(t,s)&=\int_{\mathbb R}f_{i,t}(x)e^{-wx}\,dx=P(w)M_i(t,s).
\end{aligned}\tag{EPM6}
\] The subscript \(s\) on a test denotes the sine component. All
derivatives of these tests have Gaussian tails. Their transforms are
entire, and integration by parts proves the full factor \(P(w)\). The
cosine component is even, the sine component odd, and both
\(F_i(t,0)=F_i(t,1)=0\).

Gaussian integration in the original variables gives \[
\int_{\mathbb R}e^{-x^2/(4t)}H_t(x)e^{-wx}\,dx
=2\sqrt{\pi t}\,e^{tw^2}H_0(2tw).
\tag{EPM7}
\] Fubini holds because the Gaussian dominates \(e^{|\Re w||x|}\) and
the original density has every Gaussian moment. Splitting cosine and
sine into their two exponential terms yields \[
\begin{aligned}
M_c+iM_s&=2\sqrt{\pi t}e^{tw^2-1/(4t)-iw}H_0(2tw-i)
=\frac{\sqrt{\pi t}}8e^{tw^2-1/(4t)-iw}\mathcal Q(1+itw),\\
M_c-iM_s&=2\sqrt{\pi t}e^{tw^2-1/(4t)+iw}H_0(2tw+i)
=\frac{\sqrt{\pi t}}8e^{tw^2-1/(4t)+iw}\mathcal Q(itw).
\end{aligned}\tag{EPM8}
\] The original arguments are \(\sigma_1=1+itw\), \(\sigma_0=itw\). Each
\(\mathcal Q(\sigma_j)\) means the full
\(\sigma_j(\sigma_j-1)\pi^{-\sigma_j/2}\Gamma(\sigma_j/2)\zeta(\sigma_j)\),
with EPM3--4 at exceptional points.

The two \(M_i\) never vanish together. Otherwise both products in EPM8
would vanish, placing both \(\sigma_0\) and \(\sigma_0+1\) in the open
strip \(0<\Re\sigma<1\), which is impossible. Therefore \[
\{s:M_c(t,s)=M_s(t,s)=0\}=\varnothing,\qquad
\{s:F_c(t,s)=F_s(t,s)=0\}=\{0,1\}.
\tag{EPM9}
\] At each latter point the minimum vanishing order is one. At every
original nontrivial zero at least one transform is nonzero. This holds
for every \(t>0\), without selecting a zero or time.

The source image has the direct inverse \[
H_t(x)=e^{x^2/(4t)}
[\cos(x/(2t))h_{c,t}(x)+\sin(x/(2t))h_{s,t}(x)].
\tag{EPM10}
\] It follows from \(\cos^2+\sin^2=1\), and holds on the image of the
specified map from Schwartz \(H\) to the pair. The growing-Gaussian
multiplication is not asserted continuous on arbitrary Schwartz pairs.
The filter \(T\) has the exact continuous Schwartz inverse \[
(T^{-1}f)(x)=-\int_{\mathbb R}e^{-|x-y|/2}f(y)\,dy.
\tag{EPM11}
\] The kernel has derivative jump one, so \(T(-e^{-|x|/2})=\delta_0\).
Polynomially weighted convolution estimates, moving derivatives onto the
Schwartz factor, prove continuity. The two homogeneous exponentials
\(e^{x/2},e^{-x/2}\) have no nonzero Schwartz combination. Hence this is
a two-sided inverse, returning the actual Gaussian sources when applied
to the present pair.

\subsection{2. Coupled dynamics and invertible time
map}\label{coupled-dynamics-and-invertible-time-map}

Use column vectors \(\mathbf h=(h_c,h_s)^{\mathsf T}\),
\(\mathbf M=(M_c,M_s)^{\mathsf T}\),
\(\mathbf F=(F_c,F_s)^{\mathsf T}\), and \[
J=\begin{pmatrix}0&-1\\1&0\end{pmatrix},\qquad
R(\theta)=\begin{pmatrix}\cos\theta&-\sin\theta\\\sin\theta&\cos\theta\end{pmatrix}.
\tag{EPM12}
\] Differentiation of the original integral gives
\(\partial_tH_t=-\partial_x^2H_t\). Set
\(g_\pm=e^{-x^2/(4t)\pm ix/(2t)}\), \(a_\pm=\log g_\pm\). The product
rule gives \[
\begin{aligned}
\partial_t(g_\pm H_t)&=-\partial_x^2(g_\pm H_t)
+2a_{\pm,x}\partial_x(g_\pm H_t)
+(a_{\pm,t}+a_{\pm,xx}-a_{\pm,x}^2)g_\pm H_t,\\
2a_{\pm,x}&=-x/t\pm i/t,\qquad
a_{\pm,t}+a_{\pm,xx}-a_{\pm,x}^2=1/(4t^2)-1/(2t).
\end{aligned}\tag{EPM13}
\] Taking the two real linear combinations proves \[
\partial_t\mathbf h
=-\partial_x^2\mathbf h-\frac{x}{t}\partial_x\mathbf h
+\frac1tJ\partial_x\mathbf h+
\left(\frac1{4t^2}-\frac1{2t}\right)\mathbf h .
\tag{EPM14}
\] The cross derivatives preserve the parity of the components.
Integration by parts gives \(M(Dh)=wM(h)\),
\(M(xDh)=-M(h)-w\partial_wM(h)\), hence \[
\partial_t\mathbf M
=\frac wt\partial_w\mathbf M+
\left(-w^2+\frac1{2t}+\frac1{4t^2}\right)\mathbf M
+\frac wtJ\mathbf M .
\tag{EPM15}
\] The endpoint filter contributes its full derivative: \[
\partial_t\mathbf F
=\frac wt\partial_w\mathbf F+
\left(-w^2+\frac1{2t}+\frac1{4t^2}
-\frac{2w^2}{t(w^2-1/4)}\right)\mathbf F
+\frac wtJ\mathbf F .
\tag{EPM16}
\] At \(w=\pm1/2\), use \(\mathbf F=P\mathbf M\). The apparent poles are
resolved by this exact local divisibility, without deleting the rational
term.

For \(t_0>0\), put \[
w_0=\frac t{t_0}w,\qquad
E(t,t_0,w)=\sqrt{\frac t{t_0}}
\exp\left(tw^2-t_0w_0^2-\frac1{4t}+\frac1{4t_0}\right).
\tag{EPM17}
\] The complete comparison is \[
\begin{aligned}
\mathbf M(t,\tfrac12+w)&=E(t,t_0,w)R(w_0-w)
\mathbf M(t_0,\tfrac12+w_0),\\
\mathbf F(t,\tfrac12+w)&=\frac{P(w)}{P(w_0)}E(t,t_0,w)
R(w_0-w)\mathbf F(t_0,\tfrac12+w_0).
\end{aligned}\tag{EPM18}
\] Indeed along \(tw=t_0w_0=q\), the two original products in EPM8 have
constant arguments \(iq,1+iq\). Their scalar ratio is EPM17. Recombining
the remaining \(e^{-iw},e^{iw}\) gives exactly \(R(w_0-w)\), fixing the
rotation sign.

The first map in EPM18 is an isomorphism on \(\mathcal O(\mathbb C)^2\).
Swap \(t,t_0\), substitute \(w=(t_0/t)w_0\), and reverse the rotation
for its inverse; both scalar multipliers multiply to one. The second is
an isomorphism on the specified space \(P\mathcal O(\mathbb C)^2\):
divide by \(P\) onto its image, use the first isomorphism, then multiply
by \(P\). Where \(P(w_0)=0\), use the full holomorphic quotient
\(\mathbf F(t_0)/P\), including its derivatives. This proves all
exceptional values. No isomorphism on unrestricted Schwartz pairs is
asserted. These time and argument changes act on tests; the original
zeta zeros used below remain fixed.

\subsection{3. Original prime quotient and its odd
complement}\label{original-prime-quotient-and-its-odd-complement}

The sine source is odd and is not an input to ER's even quotient. Its
exact extension is specified here. Let
\(H^{\rm all}=\mathcal S(\mathbb R)=H^{\rm ev}\oplus H^{\rm odd}\), with
projections \(\Pi_{\rm ev}h=(h(x)+h(-x))/2\) and
\(\Pi_{\rm odd}=1-\Pi_{\rm ev}\). The moment map
\(\mu(h)=(h(0),\int h)\) kills \(H^{\rm odd}\). The even functions \[
b_0=(1-2\pi x^2)e^{-\pi x^2},\quad b_1=2\pi x^2e^{-\pi x^2},
\quad \mu(b_0)=(1,0),\quad\mu(b_1)=(0,1)
\tag{EPM19}
\] split it. Let \(R_0=\mathbb C[\mathbb Q_{>0}^{\times}]\),
\(I=\ker\operatorname{aug}\), \(H_0^{\rm ev}=\ker(\mu|_{H^{\rm ev}})\).
Directly, \[
\begin{aligned}
M_0^{\rm all}&=\ker(\operatorname{aug}\otimes\mu)
=(R_0\otimes H_0^{\rm ev})\oplus(R_0\otimes H^{\rm odd})\oplus(I\otimes\mathbb C^2),\\
Q_0^{\rm all}&=M_0^{\rm all}/IM_0^{\rm all}
=H_0^{\rm ev}\oplus H^{\rm odd}\oplus(I/I^2)\otimes\mathbb C^2
=Q_0^{\rm ev}\oplus H^{\rm odd}.
\end{aligned}\tag{EPM20}
\] Here \(R_0/I=\mathbb C\) proves the first two quotient summands. The
last is \(I/I^2\). The original even quotient has a split inclusion,
with left inverse induced by \(\Pi_{\rm ev}\); the complementary
projection is induced by \(\Pi_{\rm odd}\). This defines the new domain
and its exact relation to the original one.

Prime factorization and \(t_{ab}-1=(t_a-1)+(t_b-1)+(t_a-1)(t_b-1)\)
identify \(I/I^2=\bigoplus_p\mathbb C\ell_p\) by valuations. Thus \[
\begin{aligned}
\beta[(t_p-1)\otimes f_c]
&=\ell_p\otimes\left(
H_t''(0)-\left(\frac1{2t}+\frac1{4t^2}+\frac14\right)H_t(0),
-\frac{\sqrt{\pi t}}{32}e^{-1/(4t)}\right),\\
[(t_p-1)\otimes f_s]&=0\quad\text{in }Q_0^{\rm all},\qquad
[1\otimes f_s]=f_s\ne0\quad\text{in its }H^{\rm odd}\text{ summand}.
\end{aligned}\tag{EPM21}
\] Twice differentiating EPM6 at zero proves the first entry; evaluating
EPM8 at \(s=1/2\), where both products equal one, proves the second.
Also \(1\otimes f_s\in M_0^{\rm all}\), so multiplication by \(t_p-1\)
puts it in \(IM_0^{\rm all}\). Finally \(h_s\ne0\) and injectivity of
\(T\) prove \(f_s\ne0\). These are exact quotient maps, with no new norm
assigned.

\subsection{4. Original divisor and translated
matrix}\label{original-divisor-and-translated-matrix}

Let \(i,j\in\{c,s\}\), \(\eta_c=1,\eta_s=-1\), and
\(c_\rho=\rho-\tfrac12\) for the distinct original nontrivial zeros with
multiplicities \(m_\rho\). Reflection and reality give \[
F_i(1-s)=\eta_iF_i(s),\quad F_i(\bar s)=\overline{F_i(s)},\quad
\overline{F_i(1-\bar s)}=\eta_iF_i(s).
\tag{EPM22}
\] The reflected product is \(\eta_iF_i(s)F_j(s)\), with its displayed
sign; it is not an absolute square off the critical line. For real
\(a\), let \((f_i)_a(x)=f_i(x-a)\), whose transform is
\(e^{-aw}F_i(s)\). Define \[
\begin{aligned}
B(f,g)&=\sum_\rho m_\rho\overline{M_f(1-\bar\rho)}M_g(\rho),\\
K_{ij}(a)&=B((f_i)_a,f_j)
=\sum_\rho m_\rho\eta_iF_i(\rho)F_j(\rho)e^{ac_\rho},\\
k_{ij}&=f_i^\#*f_j=\eta_i f_i*f_j,\quad f^\#(x)=\overline{f(-x)},\\
\mathcal A_{ij,a}(s)&=e^{aw}\eta_iF_i(s)F_j(s)
=M_{k_{ij}(\,\cdot+a)}(s),\quad
\mathcal A_{ij,a}(0)=\mathcal A_{ij,a}(1)=0 .
\end{aligned}\tag{EPM23}
\] All these sums converge. To give a stronger bound, for \(u\ge0\), \[
0<\Phi(u)\le D e^{9u}e^{-\pi e^{4u}},\quad
D=2\pi^2\sum_{n\ge1}n^4e^{-\pi(n^2-1)}<\infty .
\] With \(b=|\Im z|\), splitting the last exponential in two and
maximizing \((b+9)u-(\pi/2)e^{4u}\) proves \[
|H_0(z)|\le\frac{D e^{-\pi/2}}{2\pi}
\exp\left(\frac{b+9}{4}\log(b+11)\right).
\tag{EPM24}
\] Indeed \(e^{4u}\ge1+4u\) bounds the residual integral by
\(e^{-\pi/2}/(2\pi)\); elementary differentiation bounds the maximum by
\((b+9)\log(b+11)/4\). In EPM8, \(|\Im(2tw\pm i)|\le2t|\Im w|+1\),
whereas the full Gaussian modulus has exponent
\(t(\Re w)^2-t(\Im w)^2-1/(4t)\pm\Im w\). Hence, on any fixed finite
real strip \(|\Re w|\le A\), \[
|F_c(t,s)|+|F_s(t,s)|\le C_{t,A}e^{-t|\Im w|^2/2}.
\tag{EPM25}
\] The polynomial \(P\) remains in every identity. In this inequality
its polynomial magnitude times the spare Gaussian is bounded. The
\(O(|\Im w|\log(2+|\Im w|))\) exponent from EPM24 is dominated by
\(t|\Im w|^2/2\) outside a bounded interval, proving a finite constant.
The original zero count \(N(Y)=O(Y\log(Y+2))\) proves normal convergence
of EPM23 and every translation derivative on complex \(a\)-compact sets.
In particular \(K\) is entire. The original count and finite-strip
contour estimates are also used in OZG20--26, whose complete proof is
retained in the repository. No zero list is used.

Changing \(\rho\) to \(1-\rho\) and \(\bar\rho\) in the convergent
series gives \[
K_{ij}(a)\in\mathbb R\ (a\in\mathbb R),\quad
K_{ji}(a)=\eta_i\eta_jK_{ij}(a)=K_{ij}(-a),\quad
K(0)=\begin{pmatrix}q_c&0\\0&q_s\end{pmatrix}.
\tag{EPM26}
\] The off-diagonal entry is odd. Its value at zero does not remove its
other values.

Retain the original logarithmic derivatives \[
j=\zeta'/\zeta,\quad
\kappa(s)=-\tfrac12\log\pi+\tfrac12\psi_\Gamma(s/2),\quad
q(s)=1/s+1/(s-1)+\kappa(s).
\tag{EPM27}
\] The full local units EPM3--4 remain in these functions. Original
reflection gives \(j(s)+j(1-s)=-\kappa(s)-\kappa(1-s)\); the rational
derivatives of \(s(s-1)\) and their reflected partners sum to zero in
this identity.

For each finite \(N\ge1\), \(b_N=-2N-1\), set \[
\begin{aligned}
U_{N,ij}(a)&=\sum_{m=1}^N\mathcal A_{ij,a}(-2m),\\
V_{N,ij}(a)&=K_{ij}(a)+U_{N,ij}(a)-\mathcal A_{ij,a}(1),\\
\mathcal G_{N,ij}(a)&=\frac1{2\pi i}\int_{\Re s=b_N}
\mathcal A_{ij,a}(s)[\kappa(s)+\kappa(1-s)]\,ds .
\end{aligned}\tag{EPM28}
\] Every line is upward. The original argument principle, reflection,
and Mellin inversion of \(j(s)=-\sum_{n\ge2}\Lambda(n)n^{-s}\) on
\(\Re s>1\) give \[
\begin{aligned}
V_{N,ij}&=-P_{ij}(a)+\mathcal G_{N,ij},\\
\mathcal G_{N,ij}&=A_{\infty,ij}(a)+\mathcal A_{ij,a}(0)+U_{N,ij}(a),\\
K_{ij}(a)&=\mathcal A_{ij,a}(0)+\mathcal A_{ij,a}(1)
+A_{\infty,ij}(a)-P_{ij}(a).
\end{aligned}\tag{EPM29}
\] The complete expressions are \[
\begin{aligned}
P_{ij}(a)&=\sum_{n\ge2}\frac{\Lambda(n)}{\sqrt n}
[k_{ij}(a+\log n)+k_{ij}(a-\log n)],\\
A_{\infty,ij}(a)&=-(\gamma_{\rm E}+\log\pi)k_{ij}(a)\\
&\quad+\int_0^\infty
\frac{e^{-x}k_{ij}(a)-\tfrac12e^{-x/4}
[k_{ij}(a+x/2)+k_{ij}(a-x/2)]}{1-e^{-x}}\,dx .
\end{aligned}\tag{EPM30}
\] Every prime power has \(\Lambda(p^\ell)=\log p\); the two legs do not
become a factor two for a general entry and translation.

To verify the contour step, exhaust its finite width at heights
separated from zero ordinates by an inverse-polynomial distance. The
original zero count provides such a sequence. Subtracting the full \(q\)
from the factored logarithmic derivative gives polynomial growth there;
EPM25 kills the horizontal integrals. The enclosed divisor consists of
the nontrivial zeros, \(-2,\ldots,-2N\), and the pole at one, with the
signs EPM28. Moving the Gamma edge to \(1/2\) crosses residues \(-1\) at
\(0,-2,\ldots,-2N\), so the left-edge value has the additional positive
terms \(\mathcal A(0)+U_N\). The original digamma integral gives EPM30.
Gaussian convolution estimates prove absolute convergence of both prime
legs. At the Gamma origin the numerator is
\((e^{-x}-e^{-x/4})k_{ij}(a)\) minus half the weighted central second
difference, of respective orders \(O(x)\), \(O(x^2)\). Its exponential
factors and the bounded convolution give convergence at infinity. These
bounds also justify each fixed translation derivative. This proves
EPM29--30 without taking a separate infinite limit of \(U_N\).

The exact augmented comparison and inverse are \[
(K,(\mathcal A(-2m))_{m\ge1})
\longleftrightarrow
((V_N)_{N\ge0},(\mathcal A(-2m))_{m\ge1}),\quad
K=V_N-U_N+\mathcal A(1).
\tag{EPM31}
\] Both endpoints are evaluated zero for this source. Each divisor entry
and its matching Gamma boundary remain part of the map; a divergent raw
sum is not assigned a new value.

\subsection{5. All-zero matrix detection and its exact
index}\label{all-zero-matrix-detection-and-its-exact-index}

For \(\Re z>1/2\), absolute integration and summation give \[
\int_0^\infty e^{-za}K_{ij}(a)\,da
=\sum_\rho\frac{m_\rho\eta_iF_i(\rho)F_j(\rho)}{z-c_\rho}.
\tag{EPM32}
\] This is normally meromorphic on the whole plane by EPM25. Its residue
at \(c_\rho\) is \[
m_\rho
\begin{pmatrix}F_c(\rho)\\-F_s(\rho)\end{pmatrix}
\begin{pmatrix}F_c(\rho)&F_s(\rho)\end{pmatrix}.
\tag{EPM33}
\] It is nonzero of rank one by EPM9. Distinct original zeros give
distinct poles; their multiplicities are the displayed positive
integers.

If every entry of \(K(a)\) is bounded for \(a\ge0\), its Laplace
integral is holomorphic on \(\Re z>0\). It agrees with EPM32 on
\(\Re z>1/2\). The positive half-plane minus a discrete set is
connected, since any compact path can be detoured around its finitely
many intersections. The identity theorem extends equality to that
punctured half-plane. A zero with \(\Re\rho>1/2\) would contradict
holomorphy by EPM33; reflection excludes the opposite off-critical half
too.

Under RH, conversely, EPM22 at each zero is the usual complex conjugate,
and \(K\) is an absolutely convergent positive-matrix Fourier series.
Cauchy--Schwarz gives, for \(u,v\in\mathbb C^2\), \[
|u^*K(a)v|^2\le(u^*K(0)u)(v^*K(0)v)\quad(a\in\mathbb R).
\tag{EPM34}
\] Thus every entry is bounded. The same spectral Gram calculation
proves positivity of \[
\begin{pmatrix}K(0)&K(a)\\K(a)^*&K(0)\end{pmatrix}.
\tag{EPM35}
\] Conversely positivity of EPM35 for every \(a\ge0\) bounds each entry
using its two-dimensional principal restrictions. For every fixed
\(t>0\), RH is therefore equivalent to entrywise boundedness on
\([0,\infty)\), and to positivity of EPM35 for all \(a\ge0\). This
proves an equivalence, not the inequalities themselves.

Set \[
g_{i,j}=D_x^j f_i/j!,\qquad
\mathcal C_{(i,j),(\ell,k)}
=\frac{(-1)^j}{j!k!}K_{i\ell}^{(j+k)}(0).
\tag{EPM36}
\] The exact position dictionary is \(B((f_i)_a,(f_j)_b)=K_{ij}(a-b)\).
Thus EPM35 is the Gram block for positions \(0,-a\), and the kernel
\(K_{i\ell}(b-a)\) below uses source shifts \(-a,-b\). This is exactly
\(B(g_{i,j},g_{\ell,k})\): reflection contributes \(\eta_i(-1)^jw^j\).
It is Hermitian, and the entire expansion \[
K_{i\ell}(b-a)=\sum_{j,k\ge0}
\mathcal C_{(i,j),(\ell,k)}a^jb^k
\tag{EPM37}
\] converges absolutely on all compact translation pairs. Positive
semidefiniteness of every initial coefficient block yields positivity of
finite translation matrices by square Taylor truncations, including
EPM35. Under RH the spectral Gram series proves the coefficient
positivity. This is the complete coefficient criterion, with every
factorial and sign.

Its negative index is exact. Let \(\mathcal H=\ell^2(\{\rho\},m_\rho)\)
and \((\mathcal Ju)_\rho=u_{1-\bar\rho}\). Equal reflected
multiplicities make \(\mathcal J\) a self-adjoint involution. For
polynomials \(P_c,P_s\), put \[
(V_{P_c,P_s})_\rho=F_c(\rho)P_c(c_\rho)+F_s(\rho)P_s(c_\rho).
\tag{EPM38}
\] These vectors belong to \(\mathcal H\) by EPM25 and are dense.
Orthogonality to all of them makes every derivative at zero of
\(\sum_\rho m_\rho\bar u_\rho F_i(\rho)e^{zc_\rho}\) vanish, for each
\(i\). Cauchy--Schwarz and EPM25 make these series entire, hence
identically zero. Their half-line Laplace transforms and EPM32's residue
argument yield \(m_\rho\bar u_\rho F_i(\rho)=0\) for each \(\rho,i\).
Joint nonvanishing EPM9 gives \(u=0\).

With the first slot conjugate-linear,
\(\langle\mathcal J V_P,V_Q\rangle\) is EPM36. Every distinct two-point
reflection orbit supplies one negative direction of \(\mathcal J\); a
fixed critical-line point supplies none. If \(\kappa_-\) counts these
two-point orbits, allowing infinity, then \[
\sup_n n_-\left(\mathcal C|_{\{c,s\}\times\{0,\ldots,n\}}\right)=\kappa_-.
\tag{EPM39}
\] For the upper bound, projection of a negative subspace into the
negative eigenspace is injective. For the lower bound, approximate any
finite negative orthonormal family by EPM38 using density. A column
error in operator norm below \(1/4\) changes its Gram matrix by at most
\(2/4+1/16<1\), preserving negative definiteness. The finitely many
polynomials fit in one block. Multiplicities remain in the scalar
product throughout.

\subsection{6. Certified scalar and original
support}\label{certified-scalar-and-original-support}

At \(t=1/32\), AP1--28 evaluates the cosine diagonal \(q_c=K_{cc}(0)\)
from the complete arithmetic form: \[
50<q_c<64,\quad
\beta[(t_p-1)\otimes f_c]
=\ell_p\otimes\left(-J_2-\frac{1089}{4}J_0,
-\frac{\sqrt{\pi/32}}{32}e^{-8}\right),\quad
J_j=\int_0^\infty u^j e^{u^2/32}\Phi(u)\,du.
\tag{EPM40}
\] The AP proof supplies this enclosure, retaining its complete source
residual and arithmetic tails. No fresh numerical reproduction is
claimed here. The two source moments are negative. Their relation to the
positive diagonal is a relation between specified maps of the same
source. Neither the diagonal nor the moments determine EPM34--39.

For the original bounded distributive lattice \(L\) with distinct bottom
and top, retain \[
G_L(V)=\{(v,1_L):v\in V\}\cup\{z_\lambda=(0,\lambda):\lambda\ne1_L\},
\quad e=(0,1_L),\quad\tau=z_{0_L}.
\tag{EPM41}
\] Each linear map here lifts as \((v,1_L)\mapsto(Av,1_L)\),
\(z_\lambda\mapsto z_\lambda\). Substitution proves compatibility with
composition; the inverses and left inverses in EPM10--11, EPM18, EPM20
retain their stated domains. A zero amplitude under the endpoint map has
top label \(e\), whereas the absent input retains \(\tau\). The pairing
lift is \[
B^L((f,\lambda),(g,\mu))=(B(f,g),\lambda\wedge\mu).
\tag{EPM42}
\] A nonzero output requires both inputs nonzero, hence top support, so
it belongs to the carrier. Distributivity of meet over join proves the
additive identities. Independently attached lower-label function values
remain separate coordinates; they are not nonzero amplitudes at lower
points of \(G_L\). Differential operators and source transport are not
declared semiring homomorphisms.

The original signed trace, each trivial-zero entry, Gamma boundary,
prime quotient, and odd complementary source all have specified
receiving maps. All-zero detection is EPM8--9 and EPM32--33. The
remaining global sign is the actual arithmetic matrix EPM29--30 and
EPM35--39; neither a completion nor a support label has assigned it a
sign.

\subsection{7. The certified cosine test alone reaches every original
zero}\label{the-certified-cosine-test-alone-reaches-every-original-zero}

The independent complete proof
\href{ENDPOINT_RESONANT_ZERO_DETECTION.md}{ERD1--39}, read in full by
root, strengthens EPM9 on an explicit time interval. With the original
endpoint integral \[
\delta=\int_1^\infty\left(\sum_{n\ge1}e^{-\pi n^2v}\right)
(v^{-1/2}+v^{-1})\,dv,\qquad 0<\delta<2/21,
\] it proves \[
M_c(t,s)\ne0\quad(0\le\Re s\le1),\qquad
0<t\le T_*=\frac{2(\pi-1)}{\pi-1+8/21}.
\tag{EPM43}
\] To make the receiving step explicit, its full-product logarithmic
curvature is \[
\left(\frac{\zeta'}{\zeta}+\frac1s+\frac1{s-1}
-\frac12\log\pi+\frac12\psi_\Gamma(s/2)\right)'
=-\sum_\rho\frac{m_\rho}{(s-\rho)^2}.
\] The theta estimate excludes \(|\Im\rho|\le1\), and the exact endpoint
difference gives \(\sum_\rho m_\rho/[\rho(1-\rho)]=2\delta\). Taking
real parts yields \(\sum_\rho m_\rho|\Im\rho|^{-2}\le4\delta\). For
\(|v|<1\), the resulting bound on the relative phase of
\(\mathcal Q(u+iv+1)/\mathcal Q(u+iv)\) is \(4\delta|v|/(1-|v|)\). In
EPM8, \(|v|=t|\Re w|\le t/2\). Thus the relative phase between the two
cosine summands has magnitude strictly below \[
1+\frac8{21}\frac{t}{2-t}\le\pi.
\tag{EPM44}
\] They cannot cancel. ERD15--26 proves the factorization, estimates,
strict endpoint case and every original germ used in this calculation,
without a zero table.

Fix the original \(t=1/32\) and retain \(f=f_{c,1/32}\),
\(F=F_c(1/32,\cdot)\), \(K(a)=K_{cc}(a)\), \(q=K(0)\). Equations EPM6
and EPM43 prove \(F(\rho)\ne0\) at every original nontrivial zero. Thus
the scalar Laplace transform EPM32 has the nonzero residue
\(m_\rho F(\rho)^2\) at each \(c_\rho\). The proof following EPM33, now
in this one component, gives \[
\boxed{\mathrm{RH}\ \Longleftrightarrow\
K\text{ bounded on }[0,\infty)\
\Longleftrightarrow\ |K(a)|\le q\text{ for all }a\ge0,
\qquad 50<q<64.}
\tag{EPM45}
\] For the last implication towards RH, the inequality supplies
boundedness. For the converse, under RH each \(F(\rho)\) is real and the
weights \(m_\rho F(\rho)^2\) are positive; absolute convergence and the
triangle inequality yield the bound. ERD29 in fact gives
\(F(1/2+iy)=-(y^2+1/4)M_c(1/2+iy)<0\), so none of these critical-line
weights is zero either. The enclosure for \(q\) is AP's certified
arithmetic calculation, not an assumed sign.

The original translated tests satisfy, by direct expansion of the
Hermitian pairing, \[
B(f_a,f_a)=q,\qquad
Q(f_a+f)=2q+2K(a),\qquad
Q(f_a-f)=2q-2K(a),\qquad Q(g)=B(g,g).
\tag{EPM46}
\] Here the translation factors cancel exactly on each reflected
diagonal; the cross term is real by EPM26. Both endpoint moments remain
zero, with their translation multipliers retained in EPM23. Therefore
the inequalities in EPM45 are precisely positivity of these two actual
two-translate tests for every \(a\ge0\). At a hypothetical off-critical
zero, the nonzero scalar Laplace residue forces unbounded \(K\), hence a
translate and one of these two signs with negative value. This is the
contrapositive proof of the criterion, not the construction of an actual
off-critical zero.

The scalar divided-derivative matrix \[
C_{jk}=\frac{(-1)^j}{j!k!}K^{(j+k)}(0)
\tag{EPM47}
\] has the exact index \(\sup_n n_-(C|_{0,\ldots,n})=\kappa_-\). The
density proof EPM38--39 now needs only \(F\): orthogonality to all
\(F(\rho)P(c_\rho)\) makes the corresponding entire exponential series
zero, its Laplace residues give \(\bar u_\rho F(\rho)=0\), and EPM43
gives \(u=0\). The same finite-family approximation proves the index
equality. Every initial block is positive semidefinite exactly when RH
holds, by EPM37 and EPM45. The full original arithmetic evaluation of
every entry is obtained by differentiating EPM29--30 in the cosine
entry; its finite trivial-zero and Gamma coordinates remain EPM28--31.

This is the exact bridge from the newly certified scalar to the global
question: a fixed source originating in the retained endpoint residue
detects all original zeros, and its whole translated form carries their
reflection index. A positive value at \(a=0\) is one evaluated entry of
that form. The uniform bound in EPM45 remains unproved.

\end{document}
